% !TeX root = ../main.tex
\section{Introduction}
\label{sec:introduction}

Time-varying conditional quantiles describe how the lower tail, center, and upper tail of a response evolve across an ordered domain. They are useful when changes in variability or asymmetry matter as much as changes in location, as in environmental monitoring, growth assessment, and economic time series. Quantile trend filtering combines the check loss of \citet{koenker1978regression} with an $\ell_1$ penalty on discrete differences, adapting a piecewise polynomial fit to local changes in the signal \citep{tibshirani2014adaptive,madrid2022risk}.

\citet{brantley2020baseline} developed a quantile trend filtering approach for smoothly varying baseline drift, including joint non-crossing estimation, smoothing-parameter selection, and an alternating direction method of multipliers (ADMM) algorithm. Their formulation provides an effective optimization-based benchmark, especially for the smooth-baseline setting that motivates the method. At a selected regularization level, however, the optimizer supplies point estimates; posterior simulation requires a probabilistic model and an inference procedure beyond that optimization step.

We develop a Bayesian extension that retains this computational connection. Under an asymmetric Laplace working likelihood and a global Laplace prior on trend differences, the conditional posterior mode at fixed regularization parameter $\lambda$ and likelihood scale $\sigma$ is the quantile trend filtering solution with penalty $\lambda\sigma$. We infer $\lambda$ and $\sigma$ jointly with the trend instead of selecting the penalty by an external criterion. The posterior mean need not equal the optimizer, but both arise from the same check-loss and difference-penalty structure, and Section~\ref{subsec:individual} states the connection explicitly.

Related methods include Bayesian and calibrated loss-based quantile smoothing. \citet{fasiolo2021fast} developed fast calibrated additive quantile regression. For quantile trend filtering, \citet{onizuka2024fast} combined an asymmetric Laplace likelihood with shrinkage priors and developed both Gibbs sampling and calibrated variational approximations. Their horseshoe prior supports coefficient-specific shrinkage, allowing the amount of regularization to vary across the trend. The supplied Gibbs implementation uses latent-variable and local-scale updates together with a changing Gaussian precision matrix; its matrix operations motivate an alternative computational route for long signals.

Proximal Markov chain Monte Carlo provides such a route for nonsmooth models. The proximal Metropolis-adjusted Langevin algorithm (P-MALA) of \citet{pereyra2016proximal} replaces the nonsmooth negative log-density by its Moreau--Yosida envelope, whose gradient is given by a proximal map, proposes a Langevin step on the envelope, and corrects with the exact density through a Metropolis--Hastings step. Its requirement is the proximal map of the whole negative log-density. Related schemes split the density into a smooth part handled by its gradient and a nonsmooth part handled by a proximal map \citep{durmus2018efficient}, and in mean trend filtering \citet{heng2023bayesian} used proximal Bayesian computation to learn the trend, scale, and regularization parameters. Quantile trend filtering is harder because the likelihood is nonsmooth as well as the prior. The check loss has no gradient, so a splitting scheme must smooth it, which changes the target that the proposal approximates. Combined with the difference penalty it does not give a trend filtering problem whose proximal map is routinely computable. Being piecewise linear, it also carries no curvature with which to precondition a proposal.

We resolve this with the normal--exponential representation of the asymmetric Laplace distribution \citep{kozumi2011gibbs}. Conditionally on one latent variable per observation, the likelihood is Gaussian in the trend. Integrating the latent variables out returns the asymmetric Laplace likelihood exactly, so the augmentation changes the sampler and not the model. Under a global Laplace prior on the differences, the negative log full conditional of the trend is then a weighted quadratic plus the $\ell_1$ difference penalty. Its proximal map is a weighted trend filtering problem, which we solve by ADMM with a banded factorization. The augmentation is what allows the trend to be updated by P-MALA in its exact form, as \citet{pereyra2016proximal} define it, on the full conditional posterior. The Moreau--Yosida envelope is built on the entire negative log-conditional, likelihood and prior together, and its gradient costs one proximal map. The Langevin proposal is a Gaussian centred at that proximal point, so its density is explicit. The Metropolis--Hastings step is evaluated with the exact conditional density. No part of the check loss is smoothed and the target is not split, so the chain leaves the stated posterior invariant. The proximal map enters the kernel only as a deterministic function of its input, so solving it with a fixed ADMM budget affects the efficiency of the proposal, not the posterior the chain targets. The augmentation also supplies the curvature that the check loss lacks: the precision of the augmented likelihood, together with a difference term on the scale of the prior, defines the metric that preconditions the proposal.

The resulting Gibbs sampler, which we call Bayesian-PMALA, pairs this P-MALA update of the trend with conjugate updates for the latent variables, the global penalty $\lambda$ and the scale $\sigma$, and with an ancillarity--sufficiency interweaving step \citep{yu2011center} that improves the mixing of $\lambda$. The hierarchy has one global penalty per quantile rather than a vector of local shrinkage scales. Every linear system is banded, and one factorization serves both the forward and the reverse proximal maps of each move, so for a fixed ADMM iteration budget the work per sweep grows linearly with the sample size. A single default configuration, with difference order $k=2$ and fixed hyperparameters, is used for every signal, noise distribution and quantile level.

The same construction extends to several quantile levels whose curves must not cross. \citet{brantley2020baseline} impose non-crossing in the optimization; we impose it on the posterior, as an ordering constraint on the trends in a joint target formed from the single-level posteriors. Given its neighbours, the trend at each level has its single-level conditional restricted to a box. The negative log-density of that conditional is still convex, so P-MALA applies unchanged. Only the proximal map changes, and it becomes a box-constrained weighted trend filtering problem. A proposal outside the box is rejected, so the Metropolis--Hastings correction keeps its single-level form and needs no truncation constant. The constrained map is computed by the unconstrained ADMM when that solution respects the bounds, and otherwise by a primal--dual interior-point method whose Newton systems are banded, so the work per map remains linear in $n$. The latent-variable, penalty and scale updates are those of the single-level sampler, the interweaving step is restricted to the box, and every retained draw is ordered at the observed indices.

Our simulation study focuses on two established benchmarks: the varying-smoothness signal used by \citet{onizuka2024fast} and the sinusoidal smooth signal used by \citet{brantley2020baseline}, under Gaussian, Beta, and mixed-normal noise with the noise parameterization of Onizuka et al. We compare Bayesian-PMALA with the two Onizuka MCMC variants and the ADMM point estimator at quantile levels from $0.05$ to $0.95$, on identical data sets. Each competitor is most accurate on the benchmark that suits it, the horseshoe on the varying-smoothness signal and ADMM on the smooth signal, while Bayesian-PMALA is competitive on both: it attains accuracy comparable to the locally adaptive horseshoe sampler with a simpler hierarchy and banded computations, and accuracy comparable to ADMM while providing posterior inference on the trend, the penalty and the likelihood scale. A third experiment adopts the peak-detection design of \citet{brantley2020baseline}, where the target is a peak-free low-quantile baseline of a contaminated series and estimation is followed by a classification step. Because the likelihood cannot separate a burst from a movement of the trend, this setting uses a stronger prior on the penalty, fixed in advance, with the same sampler; Bayesian-PMALA then recovers the baseline far more accurately than the horseshoe sampler and comes close to ADMM in both baseline recovery and detection.

Bayesian-PMALA thus sits between computationally efficient optimization and a more elaborate locally adaptive Bayesian hierarchy. Relative to ADMM, it adds exact posterior sampling and learns the regularization and likelihood scales instead of selecting a penalty by an information criterion. Relative to the Onizuka Gibbs sampler, it replaces a vector of local shrinkage scales with one global penalty per quantile and replaces dense matrix computations with banded proximal computations; per sweep it runs $25$ to $7{,}000$ times faster than the supplied implementation for $n$ between $250$ and $4000$ (Section~\ref{sec:scalability}). Like the joint ADMM formulation, its joint version prevents crossing, here on every posterior draw.

Section~\ref{sec:methodology} presents the model, Bayesian-PMALA and its joint non-crossing extension. Appendix~\ref{app:pmala} gives the computational details, including how the unconstrained and the box-constrained proximal problems are solved. Section~\ref{sec:simulations} reports the simulation study, with additional tables in Appendix~\ref{app:simulation_results}. Section~\ref{sec:applications} presents the data applications, and Section~\ref{sec:discussion} discusses the findings and limitations.
